% !TEX program = xelatex
% ============================================================
%  Arabic Lessons Learned Template — Nibras Center
%  مشروع بحثي: تقرير الدروس المستفادة (Post-Mortem)
%  Compile with: xelatex main.tex (twice)
% ============================================================

\documentclass[11pt,a4paper]{article}

\usepackage[a4paper,margin=2.5cm,top=2.5cm,bottom=2.5cm,headheight=15pt]{geometry}
\usepackage{array}
\usepackage{booktabs}
\usepackage{tabularx}
\usepackage{enumitem}
\usepackage{float}
\usepackage{titlesec}
\usepackage{fancyhdr}
\usepackage{setspace}
\usepackage[table]{xcolor}

\usepackage{bidi}
\usepackage[no-math]{fontspec}
\usepackage[quiet]{polyglossia}

\setmainfont[Scale=1]{NotoKufiArabic-Regular.ttf}
\newfontfamily\arabicfont[Script=Arabic,Scale=0.95]{NotoKufiArabic-Regular.ttf}
\newfontfamily\englishfont[Scale=0.95]{TeX Gyre Termes}

\setdefaultlanguage[calendar=gregorian,hijricorrection=1]{arabic}
\setotherlanguage[variant=british]{english}

\usepackage[breaklinks,hidelinks]{hyperref}
\hypersetup{pdftitle={الدروس المستفادة},pdfauthor={الباحث الرئيسي}}

\titleformat{\section}{\Large\bfseries}{\thesection}{0.7em}{}
\titleformat{\subsection}{\large\bfseries}{\thesubsection}{0.6em}{}
\titlespacing*{\section}{0pt}{1.2ex plus 0.3ex}{0.8ex plus 0.2ex}

\pagestyle{fancy}
\fancyhf{}
\fancyhead[R]{\small\itshape الدروس المستفادة}
\fancyhead[L]{\small اسم المشروع}
\fancyfoot[C]{\small\thepage}
\renewcommand{\headrulewidth}{0.3pt}

\newcolumntype{L}[1]{>{\raggedright\arraybackslash}p{#1}}
\newcolumntype{C}[1]{>{\centering\arraybackslash}p{#1}}

\definecolor{goodColor}{HTML}{27AE60}
\definecolor{badColor}{HTML}{C0392B}
\definecolor{improveColor}{HTML}{D4A017}
\definecolor{headerColor}{HTML}{1A3C6B}

\setstretch{1.4}

\begin{document}

% ============================================================
% TITLE BLOCK
% ============================================================
\begin{center}
{\LARGE\bfseries تقرير الدروس المستفادة}\\[0.3em]
{\Large (\LR{Lessons Learned / Post-Mortem})}\\[0.5em]
{\large عنوان المشروع البحثي}\\[1em]
\end{center}

\begin{table}[H]
\centering
\renewcommand{\arraystretch}{1.4}
\begin{tabularx}{\textwidth}{L{4cm} X L{4cm} X}
\toprule
\textbf{اسم المشروع:} & --- & \textbf{رقم المشروع:} & --- \\
\textbf{الباحث الرئيسي:} & --- & \textbf{مدير المشروع:} & --- \\
\textbf{تاريخ البدء:} & --- & \textbf{تاريخ الانتهاء:} & --- \\
\textbf{المدة الفعلية:} & --- & \textbf{المدة المخططة:} & --- \\
\textbf{الميزانية الفعلية:} & --- & \textbf{الميزانية المخططة:} & --- \\
\textbf{تاريخ التقرير:} & \today & \textbf{مقدم التقرير:} & --- \\
\bottomrule
\end{tabularx}
\end{table}

\vspace{1em}

% ============================================================
% 1. PROJECT SUMMARY
% ============================================================
\section{ملخص المشروع}
\label{sec:summary}

% TODO: اكتب ملخصاً مختصراً للمشروع وأهدافه ونتائجه الرئيسية.
\vspace{3em}

% ============================================================
% 2. WHAT WENT WELL
% ============================================================
\section{ما سار بشكل جيد}
\label{sec:good}

\begin{table}[H]
\centering
\renewcommand{\arraystretch}{1.5}
\begin{tabularx}{\textwidth}{C{0.8cm} L{5cm} X L{3cm}}
\toprule
\textbf{رقم} & \textbf{النجاح} & \textbf{التفاصيل} & \textbf{العامل الرئيسي} \\
\midrule
1 & تعاون الفريق & تواصل ممتاز بين أعضاء الفريق & اجتماعات أسبوعية منتظمة \\
2 & جودة البيانات & دقة عالية في جمع البيانات & تدريب مسبق للفنيين \\
3 & الالتزام بالجدول & إنجاز المرحلة الأولى في الوقت المحدد & تخطيط دقيق \\
4 & قبول الورقة للنشر & قبول الورقة من أول تقديم & مراجعة داخلية صارمة \\
\bottomrule
\end{tabularx}
\end{table}

% ============================================================
% 3. WHAT WENT WRONG
% ============================================================
\section{ما لم يسر بشكل جيد}
\label{sec:bad}

\begin{table}[H]
\centering
\renewcommand{\arraystretch}{1.5}
\begin{tabularx}{\textwidth}{C{0.8cm} L{5cm} X L{3cm}}
\toprule
\textbf{رقم} & \textbf{المشكلة} & \textbf{التفاصيل} & \textbf{السبب الجذري} \\
\midrule
1 & تأخر الجهاز & تأخر تسليم جهاز القياس 3 أسابيع & تأخر المورد \\
2 & نقص العينات & صعوبة في تجنيد المشاركين & عدم كفاية الإعلان \\
3 & تجاوز الميزانية & تجاوز بند المعدات 10\% & أسعار أعلى من المتوقع \\
4 & تأخر التحليل & استغراق التحليل وقتاً أطول & نقص في الخبرة التقنية \\
\bottomrule
\end{tabularx}
\end{table}

% ============================================================
% 4. LESSONS LEARNED
% ============================================================
\section{الدروس المستفادة}
\label{sec:lessons}

\subsection{دروس إيجابية (كرّرها في المشاريع القادمة)}
\begin{enumerate}
    \item \textbf{التخطيط المبكر:} بدء التخطيط قبل الموعد الرسمي بمدة كافية يقلل من التأخيرات.
    \item \textbf{الاجتماعات المنتظمة:} الاجتماعات الأسبوعية تحافظ على تواصل الفريق وتكشف المشاكل مبكراً.
    \item \textbf{التدريب المسبق:} تدريب الفريق قبل بدء التجارب يحسن جودة البيانات.
    \item \textbf{المراجعة الداخلية:} مراجعة الورقة داخلياً قبل التقديم يزيد فرص القبول.
\end{enumerate}

\subsection{دروس سلبية (تجنّبها في المشاريع القادمة)}
\begin{enumerate}
    \item \textbf{الاعتماد على مورد واحد:} يجب البحث عن مصادر بديلة للأجهزة لتجنب التأخير.
    \item \textbf{التقليل من تجنيد المشاركين:} يجب تخصيص ميزانية كافية للإعلان وتجنيد العينات.
    \item \textbf{عدم احتساب هامش الميزانية:} يجب إضافة هامش 15\% لكل بند في الميزانية.
    \item \textbf{نقص الخبرة التقنية:} يجب تقييم المهارات المطلوبة وتدريب الفريق مسبقاً.
\end{enumerate}

% ============================================================
% 5. RECOMMENDATIONS
% ============================================================
\section{توصيات للمشاريع القادمة}
\label{sec:recommendations}

\begin{table}[H]
\centering
\renewcommand{\arraystretch}{1.5}
\begin{tabularx}{\textwidth}{C{0.8cm} L{4cm} X L{3cm}}
\toprule
\textbf{رقم} & \textbf{التوصية} & \textbf{الإجراء المقترح} & \textbf{الأولوية} \\
\midrule
1 & تخطيط مبكر للموارد & طلب الأجهزة قبل شهر من الحاجة & \color{badColor}\textbf{عالية} \\
2 & ميزانية إعلان كافية & تخصيص 5\% من الميزانية للإعلان & \color{badColor}\textbf{عالية} \\
3 & هامش ميزانية 15\% & إضافة هامش طوارئ لكل بند & \color{improveColor}\textbf{متوسطة} \\
4 & تدريب الفريق & ورشة عمل قبل بدء التجارب & \color{improveColor}\textbf{متوسطة} \\
5 & نسخ احتياطية يومية & نظام نسخ احتياطي تلقائي & \color{badColor}\textbf{عالية} \\
6 & مصادر بديلة & قائمة موردين بديلين لكل جهاز & \color{improveColor}\textbf{متوسطة} \\
\bottomrule
\end{tabularx}
\end{table}

% ============================================================
% 6. TEAM FEEDBACK
% ============================================================
\section{ملاحظات الفريق}
\label{sec:feedback}

\begin{table}[H]
\centering
\renewcommand{\arraystretch}{1.5}
\begin{tabularx}{\textwidth}{L{3cm} X}
\toprule
\textbf{العضو} & \textbf{أبرز ملاحظاته} \\
\midrule
الباحث الرئيسي & --- \\
باحث مشارك أول & --- \\
باحث مشارك ثانٍ & --- \\
فني المختبر & --- \\
مدير المشروع & --- \\
\bottomrule
\end{tabularx}
\end{table}

% ============================================================
% 7. CLOSURE
% ============================================================
\section{إغلاق المشروع}
\label{sec:closure}

\begin{table}[H]
\centering
\renewcommand{\arraystretch}{1.5}
\begin{tabularx}{\textwidth}{L{6cm} C{3cm} X}
\toprule
\textbf{بند الإغلاق} & \textbf{الحالة} & \textbf{ملاحظات} \\
\midrule
تسليم جميع المخرجات & \color{goodColor}\textbf{مكتمل} & --- \\
تسليم التقرير النهائي & \color{goodColor}\textbf{مكتمل} & --- \\
تسوية الميزانية & \color{goodColor}\textbf{مكتمل} & --- \\
أرشفة الوثائق & \color{goodColor}\textbf{مكتمل} & --- \\
تسليم المعدات & \color{goodColor}\textbf{مكتمل} & --- \\
تقييم الفريق & \color{goodColor}\textbf{مكتمل} & --- \\
\bottomrule
\end{tabularx}
\end{table}

\vspace{2em}

\begin{table}[H]
\centering
\renewcommand{\arraystretch}{2}
\begin{tabularx}{\textwidth}{L{4cm} X X}
\toprule
& \textbf{الاسم} & \textbf{التوقيع والتاريخ} \\
\midrule
الباحث الرئيسي & --- & \rule{5cm}{0.4pt} \\
مدير المشروع & --- & \rule{5cm}{0.4pt} \\
\bottomrule
\end{tabularx}
\end{table}

\end{document}
